% =============================================================================
% core/fonts.tex — famílias tipográficas
% -----------------------------------------------------------------------------
%   \cvSans            Calibri (Regular / Bold / Italic / Bold Italic)
%   \cvSerif           Noto Serif Bengali Regular  (\bfseries = SemiBold)
%   \cvSerifMedium     Noto Serif Bengali Medium   (nome no cabeçalho)
%   \cvSerifSemiBold   Noto Serif Bengali SemiBold (títulos de seção)
%   \cvBulletFont      Arial (apenas o marcador ●)
%
% Calibri é proprietária. Ordem de busca:
%   1. fonts/calibri.ttf, calibrib.ttf, calibrii.ttf, calibriz.ttf
%   2. Calibri instalada no sistema
%   3. Carlito (incluída em fonts/) — clone métrico do Calibri: mesmas
%      larguras e alturas, portanto a mesma diagramação.
%
% O Docs aplica kerning, mas não ligaduras (fi, fl...): Ligatures=NoCommon.
% =============================================================================

\RequirePackage{fontspec}

\defaultfontfeatures{Ligatures={TeX,NoCommon}, Kerning=On}

% --- Calibri -----------------------------------------------------------------
% \cv@sans{<comando>} aplica a mesma escolha a \cvSans e à fonte padrão.
\makeatletter
\chardef\cv@sanssource=0 % 0 = Carlito, 1 = fonts/calibri*.ttf, 2 = sistema
\IfFileExists{fonts/calibri.ttf}
  {\chardef\cv@sanssource=1 }
  {\IfFontExistsTF{Calibri}{\chardef\cv@sanssource=2 }{}}
\def\cv@sans#1{%
  \ifcase\cv@sanssource
    #1{Carlito}[
      Path = fonts/, Extension = .ttf,
      UprightFont = *-Regular, BoldFont = *-Bold,
      ItalicFont = *-Italic, BoldItalicFont = *-BoldItalic]%
  \or
    #1{calibri}[
      Path = fonts/, Extension = .ttf,
      UprightFont = *, BoldFont = *b, ItalicFont = *i, BoldItalicFont = *z]%
  \or
    #1{Calibri}%
  \fi}
\cv@sans{\newfontfamily\cvSans}
\cv@sans{\setmainfont}
\makeatother

% --- Noto Serif Bengali ------------------------------------------------------
\newfontfamily\cvSerif{NotoSerifBengali}[
  Path = fonts/, Extension = .ttf, Script = Latin,
  UprightFont = *-Regular, BoldFont = *-SemiBold]
\newfontfamily\cvSerifMedium{NotoSerifBengali-Medium}[
  Path = fonts/, Extension = .ttf, Script = Latin]
\newfontfamily\cvSerifSemiBold{NotoSerifBengali-SemiBold}[
  Path = fonts/, Extension = .ttf, Script = Latin]

% --- Arial (marcador) --------------------------------------------------------
\IfFontExistsTF{Arial}{%
  \newfontfamily\cvBulletFont{Arial}
}{%
  \newfontfamily\cvBulletFont{Arimo-Regular}[Path = fonts/, Extension = .ttf]
}

% Métricas verticais usadas pelo Google Docs (em unidades de "em"):
%   ascendente = espaço acima da linha de base
%   descendente = espaço abaixo (antes do line-height)
% Calibri:            (hhea.ascender + lineGap) / upm = 1988/2048 ; 512/2048
% Noto Serif Bengali: 1092/1000 ; 502/1000
\cvDefineMetrics{sans} {0.970703125}{0.25}
\cvDefineMetrics{serif}{1.092}{0.502}
